\Article{12}
Since the function $\varphi r$ is imperceptible for every perceptible value of $r$ and continually decreases as $r$ increases, it follows at once that the integral $\int rr\varphi r.\mathrm{d}r$, extended from one perceptible value to a greater one, is still imperceptible, provided that the latter value lies within the range of values on which we can perform experiments. Yet from that property alone we could by no means conclude that the integral remains imperceptible over an arbitrarily large interval of integration. The calculations of the illustrious \textsc{Laplace} are indeed stated in a way that implies such an assumption. But as long as the nature of the function $\varphi r$ is unknown, it seems wiser to refrain from every hypothetical assumption we can dispense with. Since the constant of integration in $\int rr\varphi r.\mathrm{d}r=-\psi r$ is left to our choice, let it suffice to suppose it chosen so that $\psi r=0$ for any one perceptible value of $r$ within the range of dimensions of bodies on which we can perform experiments. In this way $\psi r$ will always be imperceptible for every other such value (positive for a smaller value and negative for a larger one). Nothing here precludes its becoming perceptible for an imperceptible value of $r$. It must be added, however, that to explain the phenomena, as the distance $r$ decreases indefinitely, the value of $\psi r$ must always remain finite, or in other words $\psi0$ must be a finite quantity. More%
\newcommand{\NearTau}{\mathord{\text{\fontspec{Liberation Serif}τ}}}%
\expandafter\def\csname HistoricalNoteText6\endcsname{If we accept the explanation of the phenomena of light in the emission theory, refraction depends on the molecular attraction exerted by the particles of the transparent body on the particles of light, and the ratio of refraction depends on the value of $\psi0$, so that
\[
\frac{c\psi0}{g}=\frac{(nn-1)kk}{8\pi^3l}
\]
where $l$ denotes the length of a pendulum that makes one swing in a second, $k$ the distance light travels in a vacuum in one second, and $n$ the ratio of the sine of the angle of incidence to the sine of the angle of refraction. Thus for water $\frac{c\psi0}{g}$ is two thousand times as great as the mean distance of the sun from the earth.}%
\SourceLeaf{52}over, clearly $\frac{c\psi r}{r}$ has the same dimensions as gravity $g$, or equivalently $\frac{c\psi r}{g}$ is a length, and hence $\frac{c\psi0}{g}$ is a definite length (depending on the nature of the bodies whose attractive forces the function $fr$ describes). We may well suspect it to be immense, but in particular cases we can scarcely assign even an approximate value to it, at least without making hypothetical assumptions\HistoricalNote{6}.

\Article{13}
In exactly the same way, in the integration $\int\psi r.\mathrm{d}r=-\theta r$ we suppose the constant chosen so that $\theta r=0$ for an arbitrary value of $r$ within the range of values for which we can perform experiments. In this way $\theta r$ will be imperceptible for every such perceptible value of $r$, though it may become perceptible for an imperceptible value. Clearly $\frac{c\theta r}{g}$ expresses the area of a two-dimensional figure and hence $\frac{\theta r}{\psi r}$ a length. But $\frac{\theta0}{\psi0}$ must be a length of imperceptible magnitude, as we now show. Since $\psi r$, starting at $r=0$, continually decreases, and decreases so rapidly that it has already become imperceptible by the time $r$ has acquired a perceptible value, the value of $r$ for which $\psi r=\frac12\psi0$ must be imperceptible. Let us denote it by $\rho$. Consider the integral $\int(\psi0-\psi r)\mathrm{d}r$, which, when extended from $r=0$ to $r=R$, has the value $=R\psi0-\theta0+\theta R$. Clearly this integral is greater than the same integral extended from $r=\rho$ to $r=R$, and that integral in turn is greater than the integral $\int(\psi0-\psi\rho)\mathrm{d}r$ between the same limits. Since the last integral has the value $=(\psi0-\psi\rho)(R-\rho)=\frac12\psi0.(R-\rho)$, we have in general for every value of $R$ (greater than $\rho$)
\[
R\psi0-\theta0+\theta R>\frac12\psi0.(R-\rho)
\]
Now if $R$ is taken to denote the value of the fraction $\frac{\theta0}{\psi0}$, this relation gives
\[
\theta R>\frac12\psi0.(R-\rho)
\]
which would be absurd if $R$ were a perceptible quantity.

\SourceLeaf{53}Despite the immense magnitude of $\psi0$, then, nothing prevents $\theta0$ from being quite a modest quantity comparable to the dimensions of bodies used in experiments.

\Article{14}
It remains to examine what this property of the function $\theta$ entails for the integral (I)
\[
\int\int\frac{\mathrm{d}t.\mathrm{d}T.\cos q.\cos Q.\theta(\mathrm{d}t,\mathrm{d}T)}{(\mathrm{d}t,\mathrm{d}T)^2}
\]
This investigation must start with a simpler one. We consider a fixed point $\mu$ on either surface and work out the integral (II)
\[
\int\frac{\mathrm{d}t.\cos q.\cos Q.\theta(\mu,\mathrm{d}t)}{(\mu,\mathrm{d}t)^2}
\]
to be extended over the whole surface $t$. Here $Q$ denotes the angle between two straight lines issuing from the point $\mu$, one directed toward the element $\mathrm{d}t$ and the other fixed. And $q$ denotes the angle between two straight lines issuing from the point of the element $\mathrm{d}t$, one directed toward the point $\mu$ and the other normal to the element and directed outward.

First we observe that if the point $\mu$ is at a perceptible distance from the surface $t$, the values of all the $\theta(\mu,\mathrm{d}t)$ will be imperceptible. In this case, therefore, the whole integral (II) will be imperceptible. Thus this integral can acquire a perceptible value only insofar as the surface $t$ has parts situated at an imperceptible distance from the point $\mu$. It clearly suffices to extend the integral (II) over such parts, disregarding all those at perceptible distances.

Next we replace $\frac{\mathrm{d}t.\cos q}{(\mu,\mathrm{d}t)^2}$ by $\pm\mathrm{d}\Pi$, where $\mathrm{d}\Pi$ denotes the element on the sphere of radius $=1$ described about the center $\mu$ onto which the element $\mathrm{d}t$ is projected as seen from the point $\mu$. The upper or lower sign applies according as the element $\mathrm{d}t$ turns its outer or inner face toward the point $\mu$. The integral (II) is thus expressed as
\[
\int\pm\mathrm{d}\Pi.\cos Q.\theta(\mu,\mathrm{d}t)
\]
and evidently the value of this integral can become perceptible only insofar as the elements $\mathrm{d}\Pi$ corresponding to imperceptible distances $(\mu,\mathrm{d}t)$ fill a region of perceptible size on the sphere.

From this we readily gather that, generally speaking, our integral remains imperceptible even when the point $\mu$ lies on the surface $t$ itself. For clearly the pro\SourceLeaf{54}jections of all elements $\mathrm{d}t$ at imperceptible distances from the point $\mu$ are at an imperceptible distance from the great circle formed on the sphere by the plane tangent to the surface $t$ at the point $\mu$. Three cases must be excepted.

1) The radii of curvature of the surface $t$ at the point $\mu$ are of imperceptible magnitude.

2) The continuity of the curvature is interrupted at the point $\mu$ or within an imperceptible distance of it (cf.\ \textit{Disquiss. gen. circa superficies curvas} art.\ 3).

3) The surface $t$ has another part at an imperceptible distance from the point $\mu$, for instance if the thickness of the region $s$ is imperceptible near that point. This case may, however, be counted with the one we shall treat in the next article.

